\subsection{级数 exp、log 和 H 的 p 进上界}

本号中假设 \(p \neq 0\)。

引理 1。设 n 是一个 \(\geqslant0\) 的整数，且 \(n=n_{0}+n_{1}p+\cdots+n_{k}p^{k}\)，其中 \(0\leqslant n_{i}\leqslant p-1\)，是 n 的 p 进展开。~令 \(\mathrm{S}(n)=n_{0}+n_{1}+\cdots+n_{k}\)。则

\[
v _ {p} (n!) = \frac {n - \mathrm{S} (n)}{p - 1}.\tag{3}
\]

  \(v_{p}(n!) = \sum_{i=1}^{n} v_{p}(i)\)，且整数 \(i\)（介于 1 与 \(n\) 之间）中满足 \(v_{p}(i) \geqslant j\) 的整数个数等于 \([n/p^{j}]\)，即 \(n/p^{j}\) 的整数部分。于是

\[
v _ {p} (n!) = \sum_ {j \geqslant 0} j ([ n / p ^ {j} ] - [ n / p ^ {j + 1} ]) = \sum_ {j \geqslant 1} [ n / p ^ {j} ].
\]

由于 \([n / p^j] = \sum_{i\geqslant j}n_i p^{i - j}\)，引理得证。

引理 2。对每个整数 \(v(n) \leqslant v(n!) \leqslant (n - 1)\theta\)，有 \(v(n) \leqslant (\log n) / (\log p)\) 且 \(n \geqslant 1\)。

由引理 1，\(v(n!) = v_{p}(n!) = (n - S(n))\theta \leqslant (n - 1)\theta\)。

另一方面，\(n \geqslant p^{v(n)}\)，从而 \(v(n) \leqslant (\log n) / (\log p)\)。

令 \(\mathbf{I} = \{\mathbf{U},\mathbf{V}\}\) 为含有两个元素的集合，并令

\[
\mathrm{H} = \sum_ {r, s \geqslant 0} \mathrm{H} _ {r, s} (\mathrm{U}, \mathrm{V}) \in \hat {\mathrm{L}} _ {\mathbf {Q}} (\mathrm{I})
\]

为豪斯多夫级数（§ 6，第 4 号，定义 1）。令 \(\mathbf{Z}_{(p)}\) 为相对于素理想 \(\mathbf{Z}\) 的 \((p)\) 的局部环，令 \((e_b)_{b \in \mathbf{B}}\) 为 \(\mathrm{L}_{\mathbf{Z}_{(p)}}(\mathrm{I})\) 在 \(\mathbf{Z}\) 上的一组基（§ 2，第 11 号，定理 1）。它也是 \(\mathrm{L}_{\mathbf{Q}}(\mathrm{I})\) 在 \(\mathbf{Q}\) 上的一组基。

命题 1。令 \(r\) 和 \(s\) 为两个 \(\geqslant 0\) 的整数。若 \(H_{r,s} = \sum_{b \in B} \lambda_b e_b\)，其中 \(\lambda_b \in \mathbf{Q}\)，是 \(H_{r,s}\) 关于基 \((e_b)_{b \in B}\) 的分解，则

\[
v _ {p} (\lambda_ {b}) \geqslant - (r + s - 1) \theta \quad for all b \in \mathbf {B}.\tag{4}
\]

环 \(\mathrm{A}_{\mathbf{Z}_{(p)}}(\mathrm{I})\) 被认作由词 \(\mathbf{Z}_{(p)}\) 生成的 \(\mathrm{A}_{\mathbf{Q}}(\mathrm{I})\) 的子-\(w \in \mathrm{Mo}(\mathrm{I})\)-模。由于 \(\mathrm{L}_{\mathbf{Z}_{(p)}}(\mathrm{I})\) 是 \(\mathrm{A}_{\mathbf{Z}_{(p)}}(\mathrm{I})\) 的直因子，

\[
\mathrm{L} _ {\mathbf {Z} _ {(p)}} (\mathrm{I}) = \mathrm{A} _ {\mathbf {Z} _ {(p)}} (\mathrm{I}) \cap \mathrm{L} _ {\mathbf {Q}} (\mathrm{I}).\tag{5}
\]

令 \(f\) 为满足 \(f \leqslant (r + s - 1)\theta < f + 1\) 的整数。关系 (4) 等价于对所有 \(v_{p}(\lambda_{b}) \geqslant -f\) 有 \(b \in B\)，即 \(H_{r,s} \in p^{-f}L_{\mathbf{Z}_{(p)}}(I)\)。但由 (5)，这也等价于 \(H_{r,s} \in p^{-f}A_{\mathbf{Z}_{(p)}}(I)\)。

由 §6 第 4 号公式 (11)，只需证明：对每个整数 \(m \geqslant 1\) 以及满足下述条件的所有整数 \(r_1, \ldots, r_m, s_1, \ldots, s_m\)，

\[
r _ {1} + \dots + r _ {m} = r, \qquad s _ {1} + \dots + s _ {m} = s,\tag{6}
\]

\[
r _ {i} + s _ {i} \geqslant 1 \quad \text { for } 1 \leqslant i \leqslant m,
\]

有

\[
v _ {p} (m. r! \dots r _ {m}! s _ {1}! \dots s _ {m}!) \leqslant f.\tag{7}
\]

但是由引理 2，\(v_{p}(r_{i}!s_{i}!) \leqslant (r_{i} + s_{i} - 1)\theta\)，且 \(v_{p}(m) \leqslant v_{p}(m!) \leqslant (m - 1)\theta\)；因此 (7) 的左端有上界

\[
\theta (m - 1 + \sum_ {i = 1} ^ {m} \left(r _ {i} + s _ {i} - 1\right)) = \theta (r + s - 1);
\]

由于它是整数，故它 \(\leqslant f\)，证明完成。

